% 01_introduction.tex

The \mcp{} (\mcp{}) was introduced in late 2024 as an open
standard for connecting LLM agents to external tools, resources,
and prompt templates~\cite{mcp_spec,anthropic_mcp_announce}. Its
adoption has been rapid: every major model provider now ships an
MCP client, and a public registry lists thousands of community
servers ranging from filesystem gateways to database frontends.
The protocol's appeal is its simplicity---a single JSON-RPC
envelope with three core methods (\texttt{tools/call},
\texttt{resources/read}, \texttt{prompts/get}) plus a thin
\texttt{initialize} handshake---but that simplicity comes at a
cost. MCP distributes \emph{isolation} across eight implicit
trust boundaries (transport, session, namespace, tool, resource,
memory, cache, authentication) without mandating concrete
enforcement at any of them. As deployments move from
single-tenant desktop setups to multi-tenant server fleets, the
question of \emph{cross-tenant isolation} becomes central: can
one tenant observe another's tool calls, results, cached state,
or session metadata? The MCP specification is silent on the
question, and mainstream reference implementations inherited the
silence.

\paragraph{Why this matters.}
A multi-tenant MCP server that serves multiple users, agents, or
organisations inherits the threat model of any multi-tenant
service: each tenant's data and computation must be confined to
that tenant. The recent empirical study by Chen et al.~(2026)
finds that runtime MCP servers routinely share stdio workers,
follow symlinks across tenant roots, and key caches without
tenant identity~\cite{chen_2026_mcp_security}. Concurrent work
argues that isolation is a first-class principle for LLM-agent
safety but provides neither a concrete boundary catalogue nor
empirical measurement~\cite{jing_2026_isolation}. The gap is
empirical and architectural: no prior work has systematically
catalogued MCP isolation failures across all eight boundaries,
mapped them to automated attack patterns, and measured their
leakage under a reproducible multi-tenant harness.

\paragraph{Our approach.}
We close the gap with a five-part study:

\begin{enumerate}
  \item \textbf{An isolation-boundary catalogue} for MCP. We
        enumerate all eight boundaries and tag each with STRIDE
        rows~\cite{stride_orig} and 63 forward ticket IDs
        ($A$-$\{TRN,SES,NSP,TOL,RES,MEM,CCH,AUT\}$-$\{nnn\}$).
  \item \textbf{A 14-concept security taxonomy} with explicit
        MCP-boundary bindings, covering prompt injection,
        confused deputy, capability-based access control,
        sandboxing, side-channel leakage, cache poisoning,
        memory poisoning, namespace collision, tool squatting,
        resource traversal, request smuggling, and authentication
        bypass.
  \item \textbf{A reproducible empirical framework} that measures
        cross-tenant leakage in MCP deployments under controlled
        multi-tenant concurrency, reporting
        \emph{leakage\_rate}, \emph{time\_to\_leak},
        \emph{defense\_overhead}, and \emph{utility\_retention}
        per boundary.
  \item \textbf{An open attack library of 50 concrete attack
        classes} covering every MCP boundary, each with
        CVSS~v3.1 scoring and pytest reproducibility.
  \item \textbf{An empirical defense comparison} at the MCP layer
        (not the OS layer) quantifying leakage reduction from
        per-tenant tool registries, tenant-prefixed cache keys,
        resource-path canonicalisation, and audience-bound
        authentication tokens---individually and in composition.
\end{enumerate}

\paragraph{Top-line findings.}
Across eight baseline attacks, four research questions, and 30
iterations per cell, we find:

\begin{itemize}
  \item \textbf{RQ-1} (baseline isolation): cache attacks dominate
        the vulnerable server's leakage volume, accounting for
        85.7\% of all observed leakage ($z = 40.75$, $p < 10^{-4}$
        vs. an even split).
  \item \textbf{RQ-2} (cache dominance): the cache boundary is the
        single most-exploitable surface; cache attacks are the
        cheapest path to cross-tenant data disclosure in our
        reference vulnerable server.
  \item \textbf{RQ-3} (prompt-injection residual): even the
        hardened secure server exhibits non-trivial leakage under
        repeated prompt-injection attacks (mean leak rate
        $0.5$), confirming that prompt injection is a structural
        property of tool-using LLM agents and not a patch
        target.
  \item \textbf{RQ-4} (defense composition): no single defense
        closes the gap, but the composition of all four defenses
        reduces leak rate to zero ($t = 22.80$, $p < 10^{-5}$ on
        a paired Welch's $t$-test), validating defense-in-depth
        at the MCP layer.
\end{itemize}

\paragraph{Contributions.}
We make five contributions to the security of multi-tenant LLM
agent infrastructure:

\begin{enumerate}
  \item \textbf{An eight-boundary catalogue} for the Model Context
        Protocol, covering all eight protocol boundaries with
        STRIDE-tagged attack surfaces and 63 forward ticket IDs
        (\S\ref{sec:threat_model}). \emph{Delta vs.\ prior work:}
        Jing et al.~(2026) propose 5 generic LLM-agent
        boundaries~\cite{jing_2026_isolation}; ours is 8
        MCP-specific boundaries with concrete STRIDE rows.
  \item \textbf{A 14-concept security taxonomy with explicit
        MCP-boundary bindings} (\S\ref{sec:background}).
        \emph{Delta vs.\ prior work:} OWASP LLM Top 10
        v1.1~\cite{owasp_llm_top10_v11} is LLM-application-level,
        not protocol-level; no prior work pins security concepts
        to MCP boundaries.
  \item \textbf{A reproducible empirical framework} that measures
        cross-tenant leakage in MCP deployments under controlled
        multi-tenant concurrency (\S\ref{sec:framework}).
        \emph{Delta vs.\ prior work:} Chen et al.~(2026)
        benchmark static scanners~\cite{chen_2026_mcp_security};
        An et al.~(2026) detect but do not measure
        leakage~\cite{an_2026_flowguard}; Liu et al.~(2026) study
        tool evolution, not tenant
        isolation~\cite{liu_2026_mcpevol}.
  \item \textbf{An open attack library of 50 concrete attack
        classes} covering every MCP boundary, each with CVSS~v3.1
        scoring and pytest reproducibility (\S\ref{sec:attacks}).
        \emph{Delta vs.\ prior work:} Wang et al.~(2024),
        Torres et al.~(2026), and Lee et al.~(2026) each
        demonstrate one failure mode in isolation; no prior work
        composes them into a single library with uniform pytest
        reproducibility~\cite{wang_2024_toolcommander,
        torres_2026_memory,lee_2026_webmcp}.
  \item \textbf{An MCP-layer defense-comparison study}
        (\S\ref{sec:evaluation}). \emph{Delta vs.\ prior work:}
        Dipta et al.~(2024) study OS-level sandbox side
        channels~\cite{dipta_2024_sandbox}; ours is
        protocol-level and measures leakage reduction under
        defense composition.
\end{enumerate}

\paragraph{Why a paper now?}
The MCP specification is under active revision, and our
empirical findings and defense blueprint can inform the next
revision (\S\ref{sec:conclusion}). We release the attack library,
the measurement framework, and the analysis notebooks as open
source so that the community can reproduce and extend our
results. The artifact is described in Appendix~A and is
available at the project repository.

\paragraph{Organisation.}
\S\ref{sec:background} introduces MCP and the isolation primer;
\S\ref{sec:threat_model} presents our threat model;
\S\ref{sec:framework} describes the measurement framework;
\S\ref{sec:attacks} catalogues the attack library;
\S\ref{sec:evaluation} answers the four research questions;
\S\ref{sec:defenses} details the defense blueprint;
\S\ref{sec:related} situates our work in the literature;
\S\ref{sec:discussion} discusses limitations and ethics; and
\S\ref{sec:conclusion} concludes with spec-change
recommendations.

\paragraph{Threat-model scope.}
We deliberately scope our threat model to the protocol layer
(\S\ref{sec:background}--\S\ref{sec:threat_model}). Physical-layer
side channels~\cite{dipta_2024_sandbox}, kernel-level compromise,
and LLM-provider compromise are out of scope and disclosed
honestly in \S\ref{sec:discussion}. Our empirical claims are
therefore about \emph{what an MCP server that follows the spec
can be made to leak}, not about what a determined adversary with
host-level access can do. This scoping is consistent with the
adjacent protocol-security literature~\cite{zhu_2025_lspfuzz,
chen_2026_mcp_security}.

\paragraph{Reproducibility.}
All artefacts --- attack library, measurement framework,
manifests, raw run outputs, analysis notebooks, summary CSVs,
and rendered figures --- are open source under the MIT license.
The pinned dependency versions in \texttt{pyproject.toml} and
the seed-controlled manifests in
\texttt{experiments/manifests/} make the results reproducible
on any machine with the documented toolchain. We ran the full
RQ-1--RQ-4 suite on three independent machines (Linux x86-64,
macOS arm64, Windows x86-64) and obtained bit-identical CSVs
modulo platform-dependent timestamps.